\documentclass{article}
\usepackage{amsmath,amssymb,amsthm,amsfonts}

\newtheorem{claim}{Claim}
\newtheorem{definition}{Definition}

\title{Differential Geometry Assignment 2}
\author{Aditya Padiyar}
\date{September 2026}

\begin{document}

\maketitle

\begin{definition}
  The n-manifold $S^n(r)$ is defined as $\{v \in \mathbb{R}^{n+1}:||v||=r\}$
\end{definition}

\section{Vector fields and reparametrization}

\begin{claim}
Suppose $\gamma$ is a smooth curve on $[a,b]$, $V$ is a parallel vector field along $\gamma$ and $\phi:[c,d] \to [a,b]$ is a diffeomorphism.
Then $V \circ \phi$ is a parallel vector field along $\gamma \circ \phi$
\end{claim}

\begin{proof}
  We have $(V \circ \phi)'(t)=V'(\phi(t))\phi'(t)$. We also have $V'(\phi(t)) \in T_{\gamma(\phi(t))}M$ since $V$ is parallel to $\gamma$.
\end{proof}

\section{Tangential and normal frames}

\begin{definition}
  We define $\mathcal{Y}(M)=\{f \in C^\infty(M,\mathbb{R}^k):f(p) \in N_pM\}$ 
\end{definition}

\begin{definition}
  A tangential frame on a neighborhood $V$ of $p$ consists of $X_i \in \mathcal{X}(V)$ such that 
  $X_i(q)$ form an orthonormal basis of $T_pM$. A normal frame is defined similarly.
\end{definition}

\begin{claim}[Gram-Schmidt]
  Suppose we have $\tilde{X}_i \in C^\infty(M,\mathbb{R}^k)$ such that $\tilde{X}_i(q)$ are linearly independent. 
  Let $X'_1 = \tilde{X}_1$ and define $X_{k+1}'$ recursively by \[ X'_{k+1} = \tilde{X}_{k+1} - \sum_{i=1}^k \frac{ \langle X'_i, \tilde{X}_{k+1} \rangle } { \langle X'_i , X'_i \rangle } X'_i \]
  Now define $X_i = \frac{X'_{i}}{||X'_i||}$. Then $X_i(q)$ are orthonormal and have the same span as $\tilde{X}_i(q)$
\end{claim}

\begin{claim}
  There is a tangential frame on some neighborhood $V$
\end{claim}

\begin{proof}
  Apply the Gram-Schmidt process to $\phi_i \in \mathcal{X}(V)$ obtained from a local parametrization $\phi \in C^\infty(\Omega,V)$
\end{proof}

\begin{claim}
  There is a normal frame on some neighborhood $V$
\end{claim}

\begin{proof}
  Define $Y_i \in \mathcal{Y}(M)$ by $Y_i(q)=\hat{\pi}_q(e_i)$ and \[V=\{q:Y_i(q) \text{ are linearly independent}\}\]
  Use the determinant function to see that $V$ is open. Finally, apply Gram-Schmidt.
\end{proof}

\begin{claim}
  Suppose $N_i \in \mathcal{Y}(V)$ are a normal frame on $V$. 
   Then \[S_p(v,w)=-\sum_{i=1}^{k-n}\langle (D_vN_i)(p),w \rangle N_i(p)\]
\end{claim}

\begin{proof}
  Get $\gamma \in C^\infty((-\epsilon,\epsilon),M)$ with $\gamma(0)=p$ and $\gamma'(0)=v$. 
  We have \[\langle N_i(\gamma(t)), \pi_{\gamma(t)}w \rangle =0\]
  Differentiating the map $t \to \pi_{\gamma(t)}w$ at $0$ we get $[(D\pi)_pv]w=S_p(v,w)$ 
  and differentiating the map $t \to N_i(\gamma(t))$ at $0$ we get $(D_vN_i)(p)$. Now using 
  the rule for differentiating inner products we get the desired result.
\end{proof}

\begin{claim}
  For the manifold $S^n(r)$ we have $S_p(v,w)=-\frac{\langle v , w \rangle}{r^2}p$
\end{claim}

\begin{proof}
  Define $\gamma(t)=r\frac{p+tv}{||p+tv||}$ and $N(p)=\frac{p}{r}$. Then $(N \circ \gamma)'(0)=\frac{v}{r}$. Now use the previous claim.
\end{proof}

\section{Local Riemannian isometry}

Suppose $F \in C^\infty(M,N)$ is a local Riemannian isometry.

\begin{itemize}
  \item For $\gamma \in C^\infty(I,M)$ define $\tilde{\gamma} \in C^\infty(I,N)$ by $F \circ \gamma$
  \item For $V$ in $\mathcal{X}(\gamma)$ define $\tilde{V} \in \mathcal{X}(\tilde{\gamma})$ by $t \to DF_{\gamma(t)}(V(t))$
\end{itemize}

\begin{claim}
  We have $DF_{\gamma(t)}[\text{cov}V(t)]=\text{cov}\tilde{V}(t)$ for all $t \in I$
\end{claim}

\begin{proof}
  Choose a local parametrization $\phi \in C^\infty(\Omega,V)$ for a neighborhood $V$ of $p=\gamma(t_0)$ 
  such that $F_{|V} \in C^\infty(V,F(V))$ is a Riemannian isometry and $F(V)$ is open. Using formula (4.15) from Chapter 4 of [LeeRM] the left hand side becomes 
  \[\dot{V}^j(F_*\phi_j)(\tilde{\gamma}(t))+V^j [DF_{\gamma(t)}(\nabla_{\gamma'(t)}\phi_j)]=\dot{V}^j(F_*\phi_j)(\tilde{\gamma}(t))+V^j[\nabla_{\tilde{\gamma}'(t)}(F_*\phi_j)] \]
  We can expand the right hand side as 
  \[\text{cov}[t \to V^j(F_*\phi_j)(\tilde{\gamma}(t))](t)\]
  Applying (4.15) again we see that this is the LHS
\end{proof}

\begin{claim}
If $\gamma$ is a geodesic then so is $\tilde{\gamma}$
\end{claim}

\begin{proof}
  Substitute $V=\gamma'$ in the previous claim.
\end{proof}

\begin{claim}
We have $DF_{\gamma(t_2)}P_{t_1t_2}^\gamma=P_{t_1t_2}^{\tilde{\gamma}}DF_{\gamma(t_1)}$
\end{claim}

\begin{proof}
  Let $v \in T_{\gamma(t_1)}M$. Let $V \in \mathcal{X}(\gamma)$ be parallel to $\gamma$ with $V(t_1)=v$. 
  Applying the left hand side to $v$ we get $\tilde{V}(t_2)$. 
  We also have $\tilde{V}(t_1)=DF_{\gamma(t_1)}v$.
  Hence applying $v$ to the RHS we also get $\tilde{V}(t_2)$
\end{proof}

\end{document}
